\section{Ubicación del curso y tesis central}

Esta clase es un tema de alta densidad dentro del curso de primavera «Advanced LLM Agents» de Berkeley, impartido por Charles Sutton, de Google DeepMind. La conferencia parece cubrir dos direcciones (Code Agent y AI Security), pero el hilo conductor real es uno solo: \textbf{cómo llevar al modelo de lenguaje de «saber escribir código» a «poder completar tareas de software abiertas»}. La estructura de la exposición se despliega en tres partes: primero regresa a la evaluación y al diseño de sistemas del coding agent, después discute la adaptación de la IA en tareas de seguridad y, por último, usa el proyecto Big Sleep para explicar por qué el método basado en agentes es válido en la búsqueda real de vulnerabilidades.

\begin{importantbox}{El criterio de ingeniería de toda la clase}
La postura de Sutton no es que «entre más grandes los parámetros del modelo, mejor», sino que «el sistema pueda cerrar el ciclo». Cerrar el ciclo incluye: definición ejecutable de la tarea, resultado verificable, trayectoria iterable, e interfaz de herramientas escalable. Basta con que falte uno de estos cuatro elementos para que el resultado del agente tienda a quedarse en el nivel de demostración.
\end{importantbox}

La primera mitad de la exposición insiste varias veces en un hecho que se pasa por alto con facilidad: desarrollar un agente no es optimizar un prompt una sola vez, sino \textbf{hacer hill climbing en un espacio combinatorio enorme}. Cuando el equipo modifica continuamente la system instruction, la forma de las herramientas, el formato de retroalimentación y la proporción del post-entrenamiento, cada paso depende de la puntuación de la evaluación como guía. Si el objetivo de evaluación está mal alineado, la dirección completa de la iteración se desvía.

\begin{knowledgebox}{Por qué importa el trasfondo interdisciplinario}
El ponente relata primero su propia trayectoria (PLN $\rightarrow$ síntesis de programas $\rightarrow$ generación de código $\rightarrow$ seguridad), no como preámbulo biográfico, sino para dejar claro que: el valor de una tecnología de IA depende de qué tan a fondo se entienda el dominio del problema. Quien solo entiende el modelo, sin conocer la ingeniería de software o los procesos de seguridad, termina optimizando el sistema para que «sepa responder preguntas», no para que «sepa hacer las cosas».
\end{knowledgebox}

\begin{warningbox}{No entiendas esta clase como «más benchmarks de seguridad»}
El punto central de la conferencia no es presentar unas cuantas tablas de clasificación más, sino explicar por qué muchos benchmarks son eficaces en una etapa temprana y dejan de serlo después, y por qué los escenarios de seguridad exigen una verificación dinámica más sólida. Quien solo memorice los nombres de los conjuntos de datos se perderá la metodología de ingeniería más valiosa.
\end{warningbox}

\begin{figure}[H]
\centering
\includegraphics[width=0.88\textwidth]{frame-01-agenda.jpg}
\caption{La agenda en tres partes presentada al inicio de la conferencia: fundamentos del Code Agent, AI Security y el caso Big Sleep \protect\footnotemark}
\end{figure}
\footnotetext{Intervalo de tiempo en el video: 00:03:22--00:03:50.}

\subsection{Resumen de la sección}
Esta sección establece el sistema de referencia para leer toda la clase: no se trata de un repaso disperso de técnicas, sino de una lección de método sobre «cómo evaluar, acotar e implementar correctamente los sistemas de agentes». Todos los detalles posteriores están al servicio de esta pregunta central.
